% ==============================================================================
% SLIDES - SEMANA 04: ATUADORES INDUSTRIAIS: PNEUMÁTICA, HIDRÁULICA E INVERSORES DE FREQUÊNCIA
% ==============================================================================

% 1. CAPA
\senaicover
  {Atuadores Industriais: Pneumática, Hidráulica e Inversores de Frequência}
  {Sistemas de atuação mecânica, controle elétrico e velocidade de esteiras.}
  {Unidade Curricular: Automação Industrial}
  {Prof. André Souza}

% 2. AGENDA
\begin{senaiframe}{Visão Geral \& Agenda}{Estrutura da Aula}
  \begin{columns}[t]
    \begin{column}{0.48\textwidth}
      \begin{tcolorbox}[senaibluecard, title={1. Conceitos Chave}]
        \begin{itemize}
          \item Fundamentos teóricos e definições.
          \item Relevância para a Indústria 4.0.
        \end{itemize}
      \end{tcolorbox}
      \vspace{4pt}
      \begin{tcolorbox}[senaiambercard, title={3. Aplicação Prática}]
        \begin{itemize}
          \item Estudo de caso na Célula Smart N1.
          \item Desenvolvimento de rotinas técnicas.
        \end{itemize}
      \end{tcolorbox}
    \end{column}
    \begin{column}{0.48\textwidth}
      \begin{tcolorbox}[senairedcard, title={2. Detalhamento Técnico}]
        \begin{itemize}
          \item Ferramentas e protocolos utilizados.
          \item Boas práticas de implementação.
        \end{itemize}
      \end{tcolorbox}
      \vspace{4pt}
      \begin{tcolorbox}[senaigreencard, title={4. Exercícios \& Avaliação}]
        \begin{itemize}
          \item Desafios de fixação do conhecimento.
          \item Integração com a plataforma FactoryHub.
        \end{itemize}
      \end{tcolorbox}
    \end{column}
  \end{columns}
\end{senaiframe}

% 3. CONTEÚDO TEÓRICO
\begin{senaiframe}{Teoria • Módulo 1}{Atuadores Industriais: Pneumática, Hidráulica e Inversores de Frequência}
  \begin{columns}[t]
    \begin{column}{0.48\textwidth}
      \begin{tcolorbox}[senaibluecard, title={Fundamentos Teóricos}]
        \begin{itemize}
          \item Cilindros de dupla ação, válvulas direcionais 5/2 vias, motores trifásicos e inversores VFD.
        \end{itemize}
      \end{tcolorbox}
    \end{column}
    \begin{column}{0.48\textwidth}
      \begin{tcolorbox}[senaigreencard, title={Prática \& Laboratório}]
        \begin{itemize}
          \item Script de controle simulado de velocidade da esteira via comandos Modbus.
        \end{itemize}
      \end{tcolorbox}
    \end{column}
  \end{columns}
\end{senaiframe}
